% ============================================================
%  Báo cáo Proposal - Đồ án môn học Kỹ thuật Máy tính (CO4041)
%  Đề tài HK261-DAGD1-133: Hệ thống microgrid năng lượng mặt trời + IoT
%  Biên dịch: pdflatex report.tex  (chạy 2-3 lần cho mục lục + tham chiếu)
% ============================================================
\documentclass[oneside,final]{hcmut-report}

% Bảng và toán
\usepackage{amsmath}
\usepackage{array,longtable,multirow,tabularx}
\usepackage{enumitem}

% Sơ đồ khối (block diagram)
\usepackage{tikz}
\usetikzlibrary{positioning,arrows.meta}
\tikzset{
  blk/.style={draw, rounded corners=2pt, align=center, inner sep=4pt,
              minimum height=9mm, font=\footnotesize},
  fl/.style={-{Latex[length=2mm]}, thick},
  fl2/.style={{Latex[length=2mm]}-{Latex[length=2mm]}, thick},
}

% Hình ảnh
\usepackage{caption,float}
\usepackage[export]{adjustbox}

% Tham chiếu chéo: dùng \Cref{}
\usepackage[nameinlink]{cleveref}
\crefname{table}{bảng}{các bảng}
\Crefname{table}{Bảng}{Các bảng}
\crefname{section}{mục}{các mục}
\Crefname{section}{Mục}{Các mục}
\crefname{figure}{hình}{các hình}
\Crefname{figure}{Hình}{Các hình}

% ------------------------------------------------------------
% Cấu hình trang bìa
% ------------------------------------------------------------
\coursename{Đồ án chuyên ngành / Đồ án môn học Kỹ thuật Máy tính}
\coursecode{CO4041 - HK261}
\reporttype{BÁO CÁO ĐỀ XUẤT (PROPOSAL)}
\topiccode{HK261-DAGD1-133}
\title{Thiết kế và phát triển hệ thống microgrid sử dụng năng lượng mặt trời với chi phí tối ưu kết hợp nền tảng IoT}
\advisor{ThS. Phạm Công Thái, TS. Lê Trọng Nhân}
\reportplace{Tp. Hồ Chí Minh}
\reporttime{Tháng 09/2026}
\footertext{Đồ án môn học Kỹ thuật Máy tính - CO4041 - HK261}

\begin{document}
\coverpage%

\clearpage
\section*{Thông tin đề tài}
\addcontentsline{toc}{section}{Thông tin đề tài}

\noindent\begin{tabular}{@{}p{4.5cm}p{10.5cm}@{}}
\textbf{Mã đề tài:} & HK261-DAGD1-133 \\[2pt]
\textbf{Học phần:} & CO4041 -- Đồ án chuyên ngành / Đồ án môn học Kỹ thuật Máy tính, HK261 \\[2pt]
\textbf{GVHD:} & ThS. Phạm Công Thái, TS. Lê Trọng Nhân \\[2pt]
\textbf{Sinh viên thực hiện:} & Lê Thanh Phú (2212581), Phan Trần Nguyên Phúc (2212643), Hoàng Ngô Thiên Phúc (2212612) \\[2pt]
\textbf{Ngày cập nhật:} & 17/09/2026 \\
\end{tabular}

\section*{Tóm tắt}
\addcontentsline{toc}{section}{Tóm tắt}

Đề tài xây dựng một mô hình microgrid quy mô nhỏ dùng năng lượng mặt trời, gồm nguồn quang điện, bộ điều khiển sạc, pin lưu trữ, bộ quản lý pin (BMS), inverter, tải và lớp giám sát/điều khiển qua IoT. Báo cáo trình bày cơ sở lý thuyết đúc kết từ tài liệu tham khảo, phương pháp thiết kế BMS 12~V và mô phỏng BMS cùng tấm pin trên Proteus, thuật toán cấu hình lại dàn pin khi bức xạ không đồng đều, kiến trúc tổng thể của hệ thống, cùng kế hoạch và phân công cho toàn bộ học phần.

Mục tiêu của báo cáo là đề xuất hướng thiết kế và lộ trình triển khai để nhận góp ý của giảng viên. Thông tin pin hiện có là 12~V, 15~Ah; loại pin, tải và mã thiết bị triển khai còn cần xác định. Nhóm dự kiến tích hợp mô-đun sạc và inverter, tập trung phần tự phát triển vào BMS, đo lường, firmware, điều phối nguồn và IoT. Báo cáo trình bày phương pháp và kế hoạch; schematic Proteus, kết quả mô phỏng, số liệu đo và BOM có giá là đầu ra của các mốc tiếp theo.


\clearpage
\tableofcontents

\clearpage

\section{Giới thiệu và mục tiêu}

\subsection{Bài toán}

Đề tài yêu cầu xây dựng một hệ thống năng lượng mặt trời có lưu trữ, hỗ trợ giám sát và điều khiển qua mobile hoặc PC, với chi phí phù hợp cho mô hình năng lượng độc lập, theo mô tả và yêu cầu của đề tài học phần. Nhóm xác định ba vấn đề cần xử lý đồng thời: công suất PV phụ thuộc điều kiện môi trường; pin lưu trữ phải được giám sát và bảo vệ; các khối nguồn, tải và giao diện điều khiển phải phối hợp theo trạng thái hệ thống.

Bức xạ và nhiệt độ làm thay đổi đặc tuyến của tấm pin \cite{yaqoob}. Khi một phần dàn pin bị che, sự không đồng đều giữa các module gây tổn hao do không tương hợp và có thể tạo nhiều cực đại trên đường công suất \cite{thanh}. Do đó nhóm nắm mô hình PV và các phương pháp cấu hình dàn pin trước khi lựa chọn phương án phần cứng.

Với pin lưu trữ, điện áp toàn bộ bộ pin không đủ để xác định trạng thái từng cell. BMS phải theo dõi điện áp cell, dòng điện, nhiệt độ và trạng thái sạc, đồng thời kết hợp cân bằng cell với bảo vệ sạc/xả. Kiến trúc BMS tích hợp IoT trong \cite{kulkarni} được nhóm dùng làm tham chiếu cho luồng đo lường, xử lý và hiển thị thông tin.

\subsection{Mục tiêu của báo cáo}
\label{sec:proposal-scope}

Trong phạm vi báo cáo này, theo phạm vi giai đoạn 1 đã thống nhất với giảng viên, nhóm tập trung vào các nội dung sau:

\begin{enumerate}
    \item Đề xuất phương pháp thiết kế BMS 12~V với bốn chức năng: ước lượng SoC, giám sát nhiệt độ, cân bằng cell và mạch bảo vệ.
    \item Xác định phương án mô phỏng chức năng BMS trên Proteus trước khi thực hiện phần cứng.
    \item Xác định mô hình và quy trình mô phỏng tấm pin PV trên Proteus.
    \item Tổng hợp lý thuyết về thuật toán PV array configuration từ tài liệu tham khảo.
    \item Lập kế hoạch và phân chia nhiệm vụ cho ba thành viên.
\end{enumerate}

Theo mức độ tổng quan đã thống nhất với giảng viên, nhóm trình bày tiêu chí lựa chọn và thời điểm chốt thiết bị. Các mã linh kiện được nhắc đến làm ví dụ hoặc tham chiếu mô phỏng; việc chọn phần cứng dựa trên đầu vào đã xác định và kết quả kiểm tra tương thích ở các mốc sau.

\subsection{Mục tiêu của toàn bộ đồ án CO4041}

Nhóm phát triển mô hình có đường cấp năng lượng từ PV đến tải và pin lưu trữ; tích hợp bộ điều khiển sạc, BMS, inverter và chức năng chuyển nguồn; thu thập dữ liệu điện áp, dòng điện, công suất, nhiệt độ và SoC; xây dựng giao diện giám sát, điều khiển và lưu lịch sử. Sau khi tích hợp, nhóm thử nghiệm với các mức tải và điều kiện nguồn khác nhau, đánh giá khả năng vận hành và lập bảng chi phí theo yêu cầu đề tài.

Trong đề tài, ``chi phí tối ưu'' được hiểu là lựa chọn phương án có chi phí phù hợp trong số các phương án đáp ứng yêu cầu chức năng đã xác định. Nhóm so sánh trên cùng quy mô công suất và dung lượng, không đặt mục tiêu chứng minh một nghiệm tối ưu toàn cục.

\section{Cơ sở lý thuyết}

\subsection{Thành phần của mô hình microgrid}

Theo phạm vi đề tài \cite{specs}, mô hình gồm nguồn PV, bộ điều khiển sạc, pin lưu trữ, BMS, inverter và tải. PV cung cấp điện một chiều; bộ điều khiển sạc điều chỉnh quá trình nạp theo bộ pin; BMS theo dõi trạng thái pin và cho phép hoặc ngắt các đường sạc/xả. Inverter chuyển điện DC thành AC cho tải phù hợp. Lớp IoT thu thập dữ liệu, hiển thị và tiếp nhận lệnh của người dùng.

Bộ điều khiển sạc và BMS có vai trò riêng: nhóm cần thiết kế hoặc chọn bộ sạc có đặc tuyến phù hợp, đồng thời dùng BMS để giám sát và bảo vệ bộ pin. Chức năng ngắt quá áp của BMS không thay thế việc điều khiển điện áp và dòng sạc.

\subsection{Đặc tính tấm pin quang điện}

Đường I--V biểu diễn quan hệ dòng điện và điện áp đầu ra; đường P--V được xác định từ:
\[
P(V)=V\,I(V).
\]

Các điểm đặc trưng gồm dòng ngắn mạch \(I_{\mathrm{SC}}\), điện áp hở mạch \(V_{\mathrm{OC}}\), điện áp và dòng tại điểm công suất cực đại \(V_{\mathrm{MPP}}, I_{\mathrm{MPP}}\). Công suất cực đại là:
\[
P_{\mathrm{MPP}}=V_{\mathrm{MPP}}\,I_{\mathrm{MPP}}.
\]

Nghiên cứu \cite{yaqoob} sử dụng điều kiện tham chiếu \(G_{\mathrm{ref}}=1000\ \mathrm{W/m^2}\) và nhiệt độ cell \(T_{\mathrm{ref}}=25^\circ\mathrm{C}\). Khi giữ nhiệt độ cố định, dòng quang sinh tăng gần tỷ lệ với bức xạ. Với các tấm silicon khảo sát trong bài, tăng nhiệt độ làm điện áp hở mạch giảm và công suất cực đại giảm. Nhiệt độ cell trong mô hình không được đồng nhất với nhiệt độ không khí xung quanh.

MPPT (Maximum Power Point Tracking) là quá trình điều chỉnh điểm làm việc để tìm công suất lớn. Khi có nhiều đỉnh trên đường P--V do che bóng, một phương pháp tìm kiếm cục bộ có thể dừng ở đỉnh chưa phải lớn nhất; vấn đề này liên quan đến việc tìm hiểu cấu hình dàn pin \cite{thanh}.

\subsection{Mô hình một diode}

Theo \cite{yaqoob} và \cite{pvpmc}, mô hình gồm nguồn dòng quang sinh, diode, điện trở song song và điện trở nối tiếp. Với các tham số quy đổi ở cấp module, phương trình dòng đầu ra là:
\[
I=I_{\mathrm{ph}}
-I_0\left[\exp\!\left(\frac{V+IR_s}{nN_sV_T}\right)-1\right]
-\frac{V+IR_s}{R_{\mathrm{sh}}},
\qquad
V_T=\frac{kT_K}{q}.
\]

Trong đó \(I_{\mathrm{ph}}\) là dòng quang sinh; \(I_0\) là dòng bão hòa diode; \(R_s\) và \(R_{\mathrm{sh}}\) là điện trở nối tiếp và song song; \(n\) là hệ số lý tưởng diode; \(N_s\) là số cell nối tiếp; \(k\) là hằng số Boltzmann; \(q\) là điện tích nguyên tố; \(T_K\) là nhiệt độ tuyệt đối.

Nhóm sẽ kiểm tra cách quy đổi tham số trước khi dựng mô hình. Nếu đã dùng \(R_s,R_{\mathrm{sh}}\) ở cấp module thì không nhân thêm số cell vào hai điện trở này. Tương tự, khi gộp \(nN_sV_T\) thành một hệ số, nhóm phải tránh đưa \(N_s\) vào hai lần \cite{pvpmc}.

Trong \cite{yaqoob}, tác giả biểu diễn các thành phần dòng bằng khối toán học và nguồn phụ thuộc trong Proteus. Nhóm dự kiến sử dụng bộ tham số công bố của SM55 để kiểm tra mô hình ban đầu, sau đó cập nhật theo tấm pin được lựa chọn. Các hệ số phụ thuộc nhiệt độ sẽ được đối chiếu với công thức, bảng tham số và chiều biến thiên trong tài liệu trước khi chạy.

\subsection{Che bóng một phần và tổn hao do không tương hợp}

Các module nối tiếp có cùng dòng nhánh, còn các module song song có cùng điện áp đầu cực. Khi bức xạ không đồng đều, đặc tuyến các module khác nhau nên điểm làm việc chung có thể khiến dàn pin mất công suất so với trường hợp từng module làm việc tại điểm tối ưu \cite{thanh}.

Cell bị che có thể chịu phân cực ngược và tiêu tán công suất, gây điểm nóng. Diode bypass tạo đường dòng đi vòng qua phần được bảo vệ, giúp hạn chế tác động này nhưng không khôi phục năng lượng mặt trời đã bị che. Cấu hình lại dàn pin hướng đến giảm phần tổn hao do không tương hợp thông qua thay đổi kết nối \cite{thanh}.

\subsection{Pin lưu trữ và trạng thái sạc}
\label{sec:soc}

Dung lượng điện của pin thường biểu diễn bằng Ah. Với điện áp danh định \(V_{\mathrm{nom}}\) và dung lượng \(Q_{\mathrm{Ah}}\), năng lượng danh định được ước lượng:
\[
E_{\mathrm{nom}}\approx V_{\mathrm{nom}}\,Q_{\mathrm{Ah}}\quad [\mathrm{Wh}].
\]

Đây là giá trị danh định; năng lượng cấp được cho tải còn phụ thuộc chế độ xả và tổn hao của hệ thống. C-rate biểu diễn dòng sạc/xả theo dung lượng, chẳng hạn \(1C\) tương ứng dòng có giá trị bằng dung lượng Ah chia cho một giờ \cite{kulkarni}.

SoC (State of Charge) biểu diễn tỷ lệ dung lượng còn lại. Tra điện áp hở mạch cần bảng đặc tính phù hợp và điều kiện nghỉ; coulomb counting tích phân dòng vào/ra, cần biết trạng thái ban đầu và dung lượng tham chiếu \cite{tigauge}. Nhóm đề xuất dùng coulomb counting cho mô hình chức năng, với quy ước dòng pack dương khi xả:
\[
z_{k+1}=\operatorname{clip}_{[0,1]}
\!\left(z_k-\frac{I_{\mathrm{bat},k}\,\Delta t}{3600\,Q_{\mathrm{Ah}}}\right),
\qquad
\mathrm{SoC}_k=100\,z_k.
\]

Ở công thức này, \(\Delta t\) tính bằng giây; mô hình ban đầu giả thiết hiệu suất coulomb bằng 1. Khi dòng đổi sang chiều sạc, \(I_{\mathrm{bat}}<0\) và SoC tăng. Do sai số dòng và giá trị khởi tạo gây sai số tích lũy, nhóm bổ sung hiệu chỉnh bằng OCV khi có bảng OCV--SoC và điều kiện nghỉ phù hợp \cite{tigauge}. \pending{bảng OCV--SoC và độ chính xác trên pin thật sẽ xác định sau khi chọn cell}

\subsection{Cân bằng cell, nhiệt độ và bảo vệ}

Các cell nối tiếp có thể lệch trạng thái sạc. Cell chạm giới hạn trước sẽ hạn chế khả năng sử dụng của bộ pin. Cân bằng thụ động tạo nhánh xả qua điện trở ở cell được chọn; năng lượng dư chuyển thành nhiệt. Khi áp dụng, cần xét điều kiện cho phép cân bằng, công suất tản và ảnh hưởng đến phép đo \cite{tibalance}.

Nhóm dự kiến giám sát nhiệt độ bằng LM35 hoặc mô hình tín hiệu tương đương. Theo datasheet \cite{lm35}, hệ số chuyển đổi là \(10\ \mathrm{mV}/^\circ\mathrm{C}\):
\[
T[^\circ\mathrm{C}]=\frac{V_{\mathrm{LM35}}[\mathrm{V}]}{0.01}.
\]

BMS cần phản ứng khi điện áp, dòng hoặc nhiệt độ vượt phạm vi cho phép. Các mức giới hạn phụ thuộc cell, bộ sạc và phần tử công suất được chọn. Nhóm không áp dụng trực tiếp ngưỡng của bộ pin trong \cite{kulkarni} cho một loại pin khác.

\subsection{Giám sát và điều khiển qua IoT}

Trong \cite{kulkarni}, bộ điều khiển thu thập dữ liệu cảm biến, hiển thị tại chỗ và chuyển dữ liệu qua ESP32 đến ThingSpeak. Nhóm sử dụng kiến trúc này làm cơ sở cho luồng giám sát. Các chức năng điều khiển tải, xác nhận lệnh và lưu trạng thái sự cố là phần nhóm đề xuất thêm cho hệ thống của mình.

Dữ liệu dự kiến gồm điện áp cell và pack, dòng pin, điện áp/dòng PV, công suất, nhiệt độ, SoC, trạng thái tải, nguồn đang sử dụng và mã lỗi. Thuật toán bảo vệ sẽ chạy tại bộ điều khiển cục bộ để việc ngắt mạch không phụ thuộc vào kết nối mạng.

\section{Phương pháp thiết kế BMS 12~V trên Proteus}

\subsection{Cấu hình dự kiến và giới hạn thiết kế}

Nhóm khảo sát bộ pin LiFePO4 4S gồm bốn cell nối tiếp, điện áp danh định khoảng 12,8~V, đại diện cho hệ lưu trữ lớp 12~V. Dung lượng định hướng ban đầu khoảng 15~Ah và sẽ được xác định lại từ tải, thời gian sử dụng và ngân sách. \pending{chốt loại cell và dung lượng ở mốc M2}

Phương án 4S cho phép khảo sát đầy đủ chức năng cân bằng cell. Nếu chuyển sang ắc-quy SLA nguyên khối, nhóm vẫn giữ bài mô phỏng nhiều cell để đáp ứng yêu cầu thiết kế cân bằng.

Nhóm thiết kế schematic BMS trong Proteus, kiểm tra từng khối rồi ghép thành mô hình hoàn chỉnh. \pending{schematic Proteus của BMS}

\subsection{Các khối chức năng và giao tiếp}

Luồng xử lý:

\begin{quote}
Bộ pin $\rightarrow$ đo điện áp, dòng, nhiệt độ $\rightarrow$ bộ điều khiển $\rightarrow$ ước lượng SoC, kiểm tra lỗi và cân bằng cell $\rightarrow$ điều khiển ngắt sạc/xả, hiển thị và xuất dữ liệu.
\end{quote}

\begin{longtable}{@{}p{2.8cm} p{4.5cm} p{6.5cm}@{}}
\toprule
\textbf{Khối} & \textbf{Đầu vào} & \textbf{Chức năng và đầu ra} \\
\midrule
\endfirsthead
\toprule
\textbf{Khối} & \textbf{Đầu vào} & \textbf{Chức năng và đầu ra} \\
\midrule
\endhead
Mô hình bộ pin & Điều kiện ban đầu, nguồn sạc, tải & Điện áp các cell và dòng pack phục vụ kiểm thử \\
Đo lường & Điện áp các nút, tín hiệu dòng và nhiệt độ & Giá trị điện áp cell/pack, dòng có dấu, nhiệt độ và cờ dữ liệu hợp lệ \\
Ước lượng SoC & Dòng pack, khoảng lấy mẫu, dung lượng và SoC ban đầu & SoC ước lượng theo thời gian \\
Cân bằng cell & Điện áp cell, chế độ vận hành, nhiệt độ & Lệnh cho nhánh xả cân bằng của từng cell \\
Bảo vệ & Điện áp, dòng, nhiệt độ, lỗi cảm biến & Quyền cho phép sạc/xả và mã nguyên nhân ngắt \\
Hiển thị/giao tiếp & Dữ liệu đo, SoC, trạng thái & LCD hoặc terminal trong Proteus; dữ liệu cho khối IoT khi tích hợp \\
\bottomrule
\caption{Các khối chức năng của BMS và giao tiếp giữa chúng.}
\label{tab:bms-blocks}
\end{longtable}

MCU dùng trong mô phỏng được chọn theo thư viện và khả năng nạp chương trình của phiên bản Proteus sử dụng. \pending{chốt MCU và kiểm tra thư viện Proteus ở đầu M2}

\subsection{Phương pháp đo lường}

Với bốn cell nối tiếp, gọi \(U_i\) là điện áp nút thứ \(i\) so với cực âm chung, \(U_0=0\). Điện áp cell thứ \(i\) là:
\[
V_{\mathrm{cell},i}=U_i-U_{i-1}.
\]

Khi đọc các nút qua cầu phân áp, nhóm chọn tỷ lệ riêng cho từng kênh theo điện áp nút lớn nhất, sau đó quy đổi về \(U_i\) trước khi lấy hiệu. Dải ADC, sai số phân áp và độ phân giải được kiểm tra trước khi chọn ngưỡng cân bằng.

Khối dòng phải phân biệt chiều sạc và xả; nhóm dùng mô hình cảm biến dòng hai chiều hoặc shunt kết hợp mạch điều hòa có điểm giữa phù hợp với ADC. Khối nhiệt độ dùng LM35 nếu có mô hình khả dụng, hoặc nguồn áp tương đương đặc tuyến cảm biến để kiểm thử chức năng.

\subsection{Phương pháp ước lượng SoC}

Nhóm kiểm thử thuật toán ở Mục~\ref{sec:soc} bằng dòng đầu vào có giá trị và thời gian xác định. SoC ban đầu được đặt trong cấu hình từng ca thử, không mặc định bằng 100\%. Trước tiên nhóm kiểm tra chiều biến thiên, đơn vị thời gian và sai số so với điện tích tính trực tiếp.

Mô hình tụ điện chỉ minh họa quá trình sạc/xả, không cung cấp đường OCV--SoC của pin thật; do đó hiệu chỉnh OCV được triển khai sau khi xác định loại cell, bảng tra và điều kiện nghỉ. Kết quả ước lượng SoC của pack không thay thế cơ chế bảo vệ theo từng cell.

\subsection{Phương pháp cân bằng cell}

Nhóm dùng cân bằng thụ động: mỗi cell có nhánh điện trở nối tiếp phần tử đóng cắt, mắc song song với cell đó. Bộ điều khiển bật nhánh phù hợp khi điện áp cell vượt mức cho phép cân bằng và chênh lệch với cell thấp nhất vượt ngưỡng đặt. Điều kiện nhiệt độ, trạng thái lỗi và ngưỡng nhả được thêm để tránh đóng/cắt liên tục.

Với mô hình lý tưởng, nhóm ước lượng dòng xả và công suất điện trở qua \(I_{\mathrm{bal}}\approx V_{\mathrm{cell}}/R_{\mathrm{bal}}\) và \(P_{\mathrm{bal}}\approx V_{\mathrm{cell}}^2/R_{\mathrm{bal}}\); giá trị cụ thể tính sau khi chọn dòng cân bằng. Trong Proteus, bước đầu dùng công tắc điều khiển lý tưởng; khi thay bằng MOSFET, nhóm thiết kế mạch kích theo điện thế của từng cell.

\subsection{Phương pháp bảo vệ và điều phối trạng thái}

\begin{longtable}{@{}p{3.6cm} p{5.0cm} p{5.3cm}@{}}
\toprule
\textbf{Chức năng} & \textbf{Đại lượng kiểm tra} & \textbf{Phản ứng} \\
\midrule
\endfirsthead
\toprule
\textbf{Chức năng} & \textbf{Đại lượng kiểm tra} & \textbf{Phản ứng} \\
\midrule
\endhead
Quá áp cell (OV) & Điện áp cell lớn nhất & Ngắt quyền sạc; chỉ cho phép hoạt động khác nếu các điều kiện tương ứng vẫn hợp lệ \\
Thấp áp cell (UV) & Điện áp cell nhỏ nhất & Ngắt quyền xả \\
Quá dòng sạc/xả (OC) & Dòng pack có dấu và giới hạn theo từng chiều & Ngắt đường liên quan; ghi nhận nguyên nhân \\
Ngắn mạch (SC) & Điều kiện dòng sự cố trong mô hình & Ngắt đầu ra; khóa khôi phục cho đến khi lỗi được xử lý \\
Nhiệt độ ngoài giới hạn & Nhiệt độ và chế độ sạc/xả & Ngắt các đường không được phép; dừng cân bằng \\
Lỗi đo lường & Tín hiệu ngoài dải hoặc không hợp lệ & Chuyển sang trạng thái khóa sạc/xả \\
\bottomrule
\caption{Các chức năng bảo vệ và phản ứng của BMS.}
\label{tab:protection}
\end{longtable}

Nhóm định nghĩa các trạng thái khởi tạo, chờ, sạc, xả và lỗi. Sau khởi tạo, hệ thống chỉ cấp quyền sạc/xả khi dữ liệu đầu vào hợp lệ. Quyết định bảo vệ có ưu tiên cao hơn lệnh vận hành, cân bằng và lệnh IoT.

Trong mô hình chức năng, nhóm tách ca sạc và ca xả để kiểm tra logic. Khi tích hợp microgrid, tải có thể nhận năng lượng trực tiếp từ PV trong lúc pin nhận phần công suất dư; vì vậy không áp dụng quy tắc ``đang sạc thì cấm toàn bộ tải'' như một yêu cầu chung.

Các tham số ngưỡng tác động, ngưỡng khôi phục, thời gian xác nhận lỗi và chế độ tự phục hồi/khóa lỗi được chọn từ datasheet và yêu cầu vận hành, ghi thành bảng cấu hình dùng chung giữa mô hình và hồ sơ kiểm thử. \pending{bảng ngưỡng bảo vệ theo datasheet của pin được chọn} Bảo vệ ngắn mạch của phần cứng được đánh giá riêng ở mốc thực nghiệm vì mô phỏng logic MCU không phản ánh tốc độ cắt và khả năng chịu dòng của mạch thật.

\section{Kế hoạch mô phỏng trên Proteus}
\label{sec:simulation}

\subsection{Trình tự và phạm vi}

Nhóm tạo hai dự án Proteus riêng cho PV và BMS và triển khai song song. Sau khi mỗi mô hình đáp ứng các ca kiểm tra, nhóm khảo sát ghép nối qua khối điều khiển sạc và đường tải. Mỗi khối cần hoàn thành mô phỏng chức năng liên quan trước khi triển khai phần cứng tương ứng; khảo sát linh kiện và phát triển giao diện bằng dữ liệu giả có thể thực hiện cùng thời gian.

Với BMS, nhóm bắt đầu bằng một kênh đo và bảo vệ, sau đó mở rộng lên nhiều cell để kiểm tra cân bằng; bài thử một cell là bước kiểm tra ban đầu trước khi dựng BMS 12~V hoàn chỉnh.

\subsection{Mô phỏng tấm pin}
\label{sec:pv-simulation}

Nhóm thực hiện các bước sau theo phương pháp trong \cite{yaqoob}:

\begin{enumerate}
    \item Lập bảng tham số của một tấm pin tham khảo (ban đầu dùng SM55 trong \cite{yaqoob}, gồm thông số danh định và tham số mô hình đã công bố).
    \item Dựng các thành phần \(I_{\mathrm{ph}}, I_D, I_{\mathrm{sh}}\) bằng khối toán và nguồn phụ thuộc trong Proteus; tạo đầu vào điều chỉnh bức xạ và nhiệt độ.
    \item Quét điểm làm việc từ gần ngắn mạch đến gần hở mạch để thu các cặp \(V,I\), sau đó tính \(P=VI\).
    \item Kiểm tra \(I_{\mathrm{SC}},V_{\mathrm{OC}},V_{\mathrm{MPP}},I_{\mathrm{MPP}},P_{\mathrm{MPP}}\) tại điều kiện tham chiếu.
    \item Thay đổi riêng bức xạ và nhiệt độ để kiểm tra xu hướng và so sánh với các đường cong trong tài liệu.
\end{enumerate}

Các ca thử cho mô phỏng tấm pin:

\begin{itemize}
    \item \textbf{PV-01.} Điều kiện \(G=1000\ \mathrm{W/m^2}\), \(T=25^\circ\mathrm{C}\); thu I--V, P--V và các điểm đặc trưng; đánh giá bằng cách đối chiếu bộ thông số SM55 trong \cite{yaqoob}.
    \item \textbf{PV-02.} Điều kiện \(G=400,700,1000\ \mathrm{W/m^2}\), giữ \(T=25^\circ\mathrm{C}\); thu họ đường I--V/P--V; kiểm tra tác động của bức xạ ở cùng nhiệt độ.
    \item \textbf{PV-03.} Điều kiện \(T=25,45,60^\circ\mathrm{C}\), giữ \(G=1000\ \mathrm{W/m^2}\); thu họ đường I--V/P--V; kiểm tra tác động của nhiệt độ ở cùng bức xạ.
\end{itemize}

Nhóm ghi sai lệch của từng thông số so với tham chiếu và chỉ so sánh định lượng khi cùng điều kiện và cùng bộ tham số. Bộ tham số SM55 phục vụ kiểm tra mô hình, không phải quyết định chọn tấm pin. Nhóm hoàn thành đồ thị và dữ liệu PV-01--03 trong giai đoạn mô phỏng chức năng.

\subsection{Mô phỏng BMS}

Nhóm viết và biên dịch chương trình cho MCU mô phỏng, nạp tệp chương trình phù hợp vào mô hình Proteus, rồi kiểm tra phần đo trước khi dùng dữ liệu cho bảo vệ hoặc SoC. Sau đó nhóm thử từng chức năng độc lập và thử lỗi kết hợp. Tệp biên dịch cho Arduino Uno chỉ dùng với mô hình tương ứng; chương trình phần cứng ESP32 được biên dịch cho ESP32 sau khi cập nhật lớp giao tiếp.

Các ca thử cho mô phỏng BMS:

\begin{itemize}
    \item \textbf{BMS-01.} Đặt các điện áp cell, dòng hai chiều và nhiệt độ đã biết; quan sát giá trị đọc đúng kênh, đúng dấu, đúng đơn vị và ghi sai số.
    \item \textbf{BMS-02.} Cho dòng sạc/xả đã biết trong thời gian xác định; SoC tăng/giảm phù hợp điện tích và không vượt phạm vi 0--100\%.
    \item \textbf{BMS-03.} Tăng một cell vượt ngưỡng OV; đường sạc ngắt, mã lỗi xác định đúng cell/điều kiện.
    \item \textbf{BMS-04.} Giảm một cell dưới ngưỡng UV; đường xả ngắt.
    \item \textbf{BMS-05.} Vượt giới hạn dòng ở từng chiều; đường tương ứng ngắt, có thời điểm tác động.
    \item \textbf{BMS-06.} Tạo sự cố ngắn mạch trong mô hình đã giới hạn dòng; đầu ra bị khóa, không tự bật lại khi lỗi còn tồn tại.
    \item \textbf{BMS-07.} Đưa nhiệt độ qua ngưỡng tác động và khôi phục; cảnh báo/ngắt và phục hồi theo đúng cấu hình.
    \item \textbf{BMS-08.} Đặt các cell lệch điện áp; nhánh cân bằng chọn đúng cell, theo dõi dòng xả và độ lệch theo thời gian.
    \item \textbf{BMS-09.} Dao động quanh ngưỡng, lỗi cảm biến hoặc khởi động lại; tránh rung đóng/cắt, đầu ra khởi động ở trạng thái khóa.
    \item \textbf{BMS-10.} Có nhiều lỗi đồng thời hoặc lệnh bật tải khi đang lỗi; giữ ưu tiên bảo vệ, chỉ phục hồi khi các điều kiện cho phép đều thỏa.
\end{itemize}

Với nguồn áp lý tưởng, nhóm kiểm tra việc chọn nhánh cân bằng nhưng điện áp nguồn không tự giảm; để đánh giá diễn biến chênh lệch cell, nhóm dùng mô hình lưu trữ có động học (mô hình tụ hoặc mạch tương đương pin được xác định tham số). Kết quả từ mô hình tụ thể hiện hành vi mạch, không dùng để công bố thời gian cân bằng hay thời gian sử dụng pin thật. Nhóm hoàn thành dữ liệu và kết luận BMS-01--10 trong giai đoạn mô phỏng chức năng.

\subsection{Hồ sơ mô phỏng}

Mỗi lần chạy lưu phiên bản mô hình, tham số đầu vào, mã nguồn và tệp biên dịch MCU đối với mô hình BMS, đồ thị/dữ liệu và kết luận cho từng ca thử, kèm điều kiện và mã ca thử. Mô hình PV bằng các khối phương trình và nguồn phụ thuộc có thể chạy độc lập với MCU. Quy trình thao tác chi tiết được tách trong tài liệu hướng dẫn mô phỏng đi kèm bộ tài liệu đề tài.

Mô phỏng điện áp và tín hiệu nhiệt cho phép kiểm tra phản ứng của BMS ở mức logic. Thực nghiệm tập trung vào sai số đo, nhiệt độ tại các vị trí giám sát và phản ứng bảo vệ trong các điều kiện thử đã xác định. Nhóm sử dụng dữ liệu nhà sản xuất cho đặc tính điện hóa và tuổi thọ pin; nghiên cứu lão hóa và thử phá hủy nằm ngoài phạm vi học kỳ.

\section{Tìm hiểu thuật toán PV array configuration}
\label{sec:pv-configuration}

\subsection{Cấu hình tĩnh và cấu hình lại dàn pin}

PV array configuration mô tả cách nối các module thành dàn pin. Các cấu hình được trình bày trong \cite{thanh} gồm nối tiếp, song song, nối tiếp--song song, Total-Cross-Tied (TCT), Bridge-Link và Honey-Comb. Reconfiguration là việc thay đổi kết nối trong quá trình vận hành để thích ứng với phân bố bức xạ.

Trong cấu hình TCT được xét ở \cite{thanh}, các module trong một hàng mắc song song và các hàng liên kết nối tiếp. Khi tổng khả năng cấp dòng của các hàng không đồng đều, điểm làm việc của dàn bị ảnh hưởng. Ý tưởng của reconfiguration là phân bố lại các module giữa các hàng để giảm sự chênh lệch này.

Diode bypass và reconfiguration có vai trò khác nhau. Bypass tạo đường dòng đi vòng qua phần được bảo vệ; reconfiguration thay đổi cách ghép các module. Việc chia nhóm module và bố trí bypass có thể được cân nhắc khi làm phần cứng, nhưng không thay thế nội dung tìm hiểu thuật toán trong báo cáo.

\subsection{Cân bằng bức xạ và chỉ số EI}

Theo \cite{thanh}, hệ thống sử dụng thông tin điện áp, dòng điện cùng mô hình module để ước lượng bức xạ. Ký hiệu \(G_{ij}\) là bức xạ của module thứ \(j\) trong hàng \(i\), \(n_i\) là số module ở hàng đó và \(m\) là số hàng. Tổng bức xạ quy ước của hàng là:
\[
G_i=\sum_{j=1}^{n_i}G_{ij}.
\]

Với số hàng cố định, mức phân bố đều tương ứng:
\[
\overline{G}=\frac{1}{m}\sum_{i=1}^{m}G_i.
\]

Chỉ số cân bằng được định nghĩa:
\[
EI=\max_i(G_i)-\min_i(G_i).
\]

Các tổng \(G_i\) phục vụ so sánh khả năng cấp dòng giữa các hàng theo giả thiết mô hình, không phải phép đo năng lượng của hàng. Khi EI giảm, mức chênh lệch giữa các hàng giảm. Nghiên cứu \cite{thanh} sử dụng mục tiêu này để tìm cấu hình phù hợp, đồng thời ràng buộc số hàng theo dải điện áp đầu vào inverter.

\subsection{Thuật toán lai DP--SC}

Tài liệu \cite{thanh} mô tả cách kết hợp Dynamic Programming (DP, quy hoạch động) và Smart Choice (SC). DP được dùng cho bài toán chọn tập con module sao cho tổng bức xạ của các hàng gần cân bằng, liên quan đến bài toán Subset Sum. SC bổ sung cách lựa chọn cho những trường hợp mà phương pháp DP trong nghiên cứu chưa cho phân bố tốt.

Có thể diễn giải luồng xử lý tổng quát của phương pháp trong \cite{thanh} như sau:

\begin{enumerate}
    \item Thu nhận hoặc ước lượng bức xạ của từng module và ghi nhận cấu hình hiện tại.
    \item Xác định số hàng cùng các điều kiện kết nối cho phép.
    \item Dùng phương pháp lai DP--SC để phân nhóm module theo mục tiêu cân bằng bức xạ.
    \item Tính EI và lựa chọn cấu hình theo tiêu chí của thuật toán.
    \item Chuyển kết quả phân nhóm thành lệnh cho ma trận DES.
    \item Khảo sát đặc tuyến và công suất sau khi đổi cấu hình.
\end{enumerate}

Đây là phần diễn giải lý thuyết của thuật toán trong tài liệu, làm cơ sở khi nhóm cân nhắc mở rộng ở giai đoạn sau.

\subsection{Ma trận chuyển mạch DES}

Dynamic Electrical Scheme (DES) là phần mạch thực hiện thay đổi kết nối. Đầu ra thuật toán xác định module thuộc hàng nào; bộ điều khiển dùng thông tin này để điều khiển các công tắc tương ứng. Trong mô hình chín module của \cite{thanh}, ma trận cho phép phân bố module vào tối đa ba hàng.

Khi xét khả năng làm phần cứng, nhóm cần đánh giá số lượng công tắc, tổn hao dẫn, trình tự chuyển mạch và chi phí điều khiển. EI là chỉ tiêu trong mô hình khảo sát, chưa đủ để khẳng định mọi cấu hình EI thấp hơn đều có lợi hơn khi tính cả tổn hao và thời gian chuyển mạch.

\subsection{Phạm vi áp dụng vào đồ án}

Nhóm sử dụng \cite{thanh} để giải thích cơ chế giảm tổn hao do không tương hợp, các thành phần của hệ reconfiguration và mối liên hệ với điều kiện bức xạ. Các kết quả mô phỏng trong bài thuộc nghiên cứu của tác giả và được nhóm dùng làm tham chiếu lý thuyết.

Trong đề tài, phần này dừng ở mức tìm hiểu lý thuyết; mô phỏng hoặc chế tạo DES là hướng mở rộng tùy chọn, xem xét sau khi hoàn thành PV, BMS và hệ thống chính. MATLAB--Simulink là môi trường của nghiên cứu \cite{thanh}; các mô phỏng của nhóm đều dùng Proteus.

\section{Kiến trúc và phương án triển khai toàn hệ thống}

\subsection{Kiến trúc năng lượng và dữ liệu}

Nhóm đề xuất các đường năng lượng như Hình~\ref{fig:energy}: PV cấp qua bộ điều khiển sạc lên bus DC được quản lý, từ đó ra tải DC hoặc qua inverter cấp tải AC; pin lưu trữ trao đổi công suất hai chiều với bus qua đường sạc/xả do BMS quản lý; nguồn DC phụ trợ vào bus qua khối chọn nguồn và chống cấp ngược.

\begin{figure}[H]
\centering
\begin{adjustbox}{max width=\linewidth}
\begin{tikzpicture}[node distance=7mm and 11mm]
  \node[blk] (pv) {PV};
  \node[blk, right=of pv] (cc) {Bộ điều khiển\\sạc};
  \node[blk, right=of cc] (bus) {Bus DC\\(được quản lý)};
  \node[blk, right=of bus] (inv) {Inverter};
  \node[blk, right=of inv] (ac) {Tải AC};
  \node[blk, above=of bus] (dc) {Tải DC};
  \node[blk, below=of cc] (pin) {Pin lưu trữ};
  \node[blk, below=of bus] (bms) {BMS};
  \node[blk, below=of inv] (sel) {Chọn nguồn /\\chống cấp ngược};
  \node[blk, below=of ac] (aux) {Nguồn DC\\phụ trợ};

  \draw[fl] (pv) -- (cc);
  \draw[fl] (cc) -- (bus);
  \draw[fl] (bus) -- (inv);
  \draw[fl] (inv) -- (ac);
  \draw[fl] (bus) -- (dc);
  \draw[fl2] (pin) -- (bms);
  \draw[fl2] (bms) -- (bus);
  \draw[fl] (aux) -- (sel);
  \draw[fl] (sel) -- (bus);
\end{tikzpicture}
\end{adjustbox}
\caption{Sơ đồ khối luồng năng lượng của hệ thống.}
\label{fig:energy}
\end{figure}

Nguồn phụ trợ dự kiến là nguồn DC đã cách ly hoặc nguồn phòng thí nghiệm phù hợp. Cách ghép nguồn phụ vào bus hay đầu vào bộ sạc sẽ được xác định khi chọn bộ điều khiển nguồn. Mô hình dự kiến vận hành độc lập; chức năng chuyển nguồn không đồng nghĩa với hòa lưới.

Luồng dữ liệu và điều khiển như Hình~\ref{fig:data}: cảm biến đưa số đo về bộ điều khiển cục bộ/BMS, qua khối kết nối IoT lên dịch vụ lưu trữ và hiển thị trên giao diện mobile/PC; lệnh người dùng đi ngược lại, được bộ điều khiển kiểm tra điều kiện trước khi tác động cơ cấu chấp hành và phản hồi trạng thái thực tế.

\begin{figure}[H]
\centering
\begin{adjustbox}{max width=\linewidth}
\begin{tikzpicture}[node distance=8mm and 10mm]
  \node[blk] (sen) {Cảm biến};
  \node[blk, right=of sen] (ctrl) {Bộ điều khiển\\cục bộ / BMS};
  \node[blk, right=of ctrl] (iot) {Khối kết nối\\IoT};
  \node[blk, right=of iot] (cloud) {Dịch vụ\\lưu trữ};
  \node[blk, right=of cloud] (ui) {Giao diện\\mobile / PC};
  \node[blk, below=of ctrl] (act) {Cơ cấu\\chấp hành};

  \draw[fl] (sen) -- (ctrl);
  \draw[fl] (ctrl) -- (iot);
  \draw[fl] (iot) -- (cloud);
  \draw[fl] (cloud) -- (ui);
  \draw[fl] (ctrl) -- (act);
  \draw[fl] (ui.north) to[out=100,in=80]
        node[above, font=\scriptsize] {lệnh người dùng} (ctrl.north);
  \draw[fl] (act.south) -- ++(0,-4mm) coordinate (fb) -| (sen.south);
  \node[font=\scriptsize, below=0pt of fb] {phản hồi trạng thái};
\end{tikzpicture}
\end{adjustbox}
\caption{Sơ đồ khối luồng dữ liệu và điều khiển của hệ thống.}
\label{fig:data}
\end{figure}

Hai sơ đồ khối trên mô tả chức năng ở mức tổng quan của hệ thống, chưa bao gồm mức nối chân và bố trí phần tử công suất. \pending{sơ đồ nối chân, vị trí phần tử công suất và triển khai mạch thực tế (mốc M4--M5)}

\subsection{Phương án cho từng thành phần}

Phương án dự kiến cho từng thành phần của hệ thống:

\begin{itemize}
    \item \textbf{PV.} Chọn quy mô công suất theo tải; kiểm tra dải điện áp và dòng phù hợp bộ điều khiển sạc.
    \item \textbf{Bộ điều khiển sạc.} So sánh phương án PWM/MPPT theo yêu cầu công suất và chi phí; xác định đặc tuyến sạc theo loại pin.
    \item \textbf{Pin và BMS.} Chốt cấu hình cell, dung lượng, bảo vệ và cân bằng; hoàn thành mô phỏng trước thử nghiệm.
    \item \textbf{Inverter.} Chọn mô-đun hoặc phương án triển khai phù hợp bus và tải; kiểm tra công suất liên tục, khởi động và hiệu suất.
    \item \textbf{Chuyển nguồn.} Xác định điều kiện chọn nguồn, ngắt trước khi đóng nguồn khác khi cần, tránh cấp ngược và ghi nhận trạng thái.
    \item \textbf{Tải.} Chọn tải DC và tải AC đại diện, có các mức công suất để kiểm tra thay đổi tải.
    \item \textbf{Đo lường.} Đo dòng/áp tại PV, pin và tải ở mức cần thiết cho điều khiển và đánh giá.
    \item \textbf{IoT/giao diện.} Hiển thị hiện tại, đồ thị lịch sử, năng lượng, trạng thái nguồn/tải, cảnh báo và lệnh điều khiển.
\end{itemize}

Nhóm ưu tiên tích hợp các mô-đun công suất phù hợp nếu việc tự thiết kế toàn bộ bộ sạc hoặc inverter vượt quá nguồn lực. Công việc thiết kế và đánh giá BMS, đo lường, điều phối nguồn và IoT vẫn được giữ trong phạm vi đề tài. Quyết định dùng mô-đun hay tự thiết kế sẽ được ghi rõ trong báo cáo và BOM.

\subsection{Xác định quy mô và chi phí}

Trước khi chọn phần cứng, nhóm sẽ lập danh sách tải gồm công suất và thời gian sử dụng. Nhu cầu năng lượng tải được ước lượng:
\[
E_{\mathrm{load}}=\sum_j P_j\,t_j.
\]

Nhóm sẽ từ nhu cầu này xác định dung lượng lưu trữ, công suất PV và công suất các khối chuyển đổi, có xét tổn hao, mức dung lượng pin cho phép sử dụng và dự phòng. Các giá trị chưa có dữ liệu hiện vẫn để mở.

BOM dự kiến gồm PV, cell, bộ sạc, inverter, BMS, vi điều khiển, cảm biến, thiết bị chuyển mạch, dây nối, đầu nối, phần cơ khí và chi phí gia công. Khi so sánh phương án, nhóm sẽ ghi nguồn giá, ngày lấy giá, phần có sẵn và phần phải mua; không so trực tiếp hai phương án khác công suất hoặc thiếu chức năng.

\subsection{Kế hoạch IoT và điều khiển}

Nhóm dự kiến đánh giá ESP32 cho kết nối mạng và ThingSpeak theo kiến trúc trong \cite{kulkarni}. Nền tảng cuối cùng sẽ được chọn sau khi kiểm tra khả năng lưu dữ liệu, hiển thị và nhận lệnh. Giao diện cần:

\begin{itemize}
    \item Hiển thị điện áp, dòng, công suất, năng lượng, nhiệt độ, SoC và trạng thái tải/nguồn.
    \item Tra cứu lịch sử và thời điểm xảy ra lỗi.
    \item Gửi lệnh bật/tắt tải hoặc chọn chế độ trong phạm vi cho phép.
    \item Hiển thị kết quả thực thi và trạng thái kết nối.
\end{itemize}

Mỗi bản ghi sẽ có mốc thời gian và trạng thái chất lượng dữ liệu. Khi mất kết nối, bộ điều khiển tiếp tục bảo vệ cục bộ; giao diện phải phân biệt dữ liệu mới với dữ liệu cũ. Lệnh bật tải không được vô hiệu hóa điều kiện ngắt của BMS. Nhóm sẽ thử giao tiếp bằng terminal trên Proteus trước, sau đó kiểm tra đường truyền và ứng dụng trên thiết bị.

\section{Kế hoạch thực hiện và phân chia nhiệm vụ}

\subsection{Tiến độ dự kiến cho toàn project}

Nhóm đề xuất khung 12 tuần tính từ thời điểm bắt đầu đồ án, chia thành bảy mốc từ M1 đến M7. Mốc M1 tương ứng nội dung của báo cáo này; các mốc còn lại được điều chỉnh theo tiến độ thực tế và phản hồi của giảng viên.

Tiến độ dự kiến cho toàn bộ đồ án CO4041:

\begin{itemize}
    \item \textbf{M1 --- Proposal (tuần 1--2).} Công việc: tổng hợp tài liệu, xác định phương pháp BMS/PV, tìm hiểu reconfiguration, lập kế hoạch. Đầu ra: báo cáo có đủ nội dung yêu cầu, nguồn trích dẫn và phân công.
    \item \textbf{M2 --- Mô hình và đầu vào (tuần 3--4).} Công việc: kiểm tra Proteus, dựng PV, kiểm tra đo áp, dòng, nhiệt, xác định tải, pin và bảng tham số dự kiến. Đầu ra: PV có dữ liệu đối chiếu, khối đo đọc đúng tín hiệu, bộ thông số đủ để lập mô hình BMS.
    \item \textbf{M3 --- BMS hoàn chỉnh trong mô phỏng (tuần 5--6).} Công việc: mở rộng nhiều cell, thêm SoC, cân bằng, bảo vệ và thử lỗi kết hợp. Đầu ra: dự án Proteus, chương trình, cấu hình và kết quả PV-01--03, BMS-01--10; xử lý lỗi ảnh hưởng chức năng.
    \item \textbf{M4 --- Thiết kế tích hợp (tuần 7).} Công việc: chọn bộ sạc, inverter, tải và chuyển nguồn, khảo sát ghép khối trên Proteus, chốt BOM và giao tiếp. Đầu ra: sơ đồ kết nối, dải điện áp/dòng phù hợp, danh mục vật tư và kế hoạch thử phần cứng.
    \item \textbf{M5 --- Phần cứng và IoT (tuần 8--9).} Công việc: triển khai mạch đo/BMS, tích hợp từng khối nguồn ở mức phù hợp, phát triển firmware và giao diện. Đầu ra: các khối chạy riêng, hiệu chuẩn ban đầu, dữ liệu và lệnh được truyền có phản hồi.
    \item \textbf{M6 --- Tích hợp và thực nghiệm (tuần 10--11).} Công việc: ghép hệ thống, thử tải, điều kiện nguồn, chuyển nguồn, mất mạng và bảo vệ, thu số liệu chi phí/hiệu suất. Đầu ra: nhật ký thử nghiệm và dữ liệu, danh sách lỗi đã xử lý hoặc giới hạn còn lại.
    \item \textbf{M7 --- Hoàn thiện và bàn giao (tuần 12).} Công việc: dự phòng sửa lỗi, kiểm tra lại ca bị ảnh hưởng, hoàn thiện báo cáo và hồ sơ. Đầu ra: mô hình minh họa, mã nguồn, schematic, kết quả, BOM và báo cáo cuối theo yêu cầu học phần.
\end{itemize}

Việc mua linh kiện chỉ triển khai sau khi đã xác định thông số và khả năng tương thích của khối liên quan. Khảo sát nhà cung cấp có thể làm trong M2--M3 để ước lượng thời gian cung ứng. Phần giao diện có thể phát triển với dữ liệu giả lập ở M5, nhưng đánh giá IoT hoàn chỉnh phải dùng dữ liệu từ thiết bị ở M6.

\subsection{Phân chia nhiệm vụ}

Phân công dưới đây kế thừa hướng phụ trách trong bản kế hoạch ban đầu và mở rộng cho toàn bộ project.

Phân chia nhiệm vụ giữa ba thành viên:

\begin{itemize}
    \item \textbf{Lê Thanh Phú (2212581).} Tại mốc proposal: tổng hợp BMS, SoC, nhiệt độ, cân bằng và bảo vệ, xây dựng phương pháp mô phỏng BMS. Sau proposal: dựng BMS trong Proteus, phát triển logic cục bộ, triển khai BMS phần cứng, phối hợp thử lỗi. Sản phẩm phụ trách: mô hình BMS, cấu hình bảo vệ, chương trình và hồ sơ kiểm thử BMS.
    \item \textbf{Phan Trần Nguyên Phúc (2212643).} Tại mốc proposal: tổng hợp PV, shading, TCT/DES và thuật toán DP--SC, lập kế hoạch mô phỏng PV. Sau proposal: dựng PV, chủ trì bộ sạc, inverter, tải và chuyển nguồn, theo dõi BOM khối công suất. Sản phẩm phụ trách: mô hình PV, dữ liệu I--V/P--V, phương án cấp nguồn và kết quả thử năng lượng.
    \item \textbf{Hoàng Ngô Thiên Phúc (2212612).} Tại mốc proposal: tổng hợp đo lường, kiến trúc IoT, kế hoạch và biên tập báo cáo. Sau proposal: chủ trì cảm biến/hiệu chuẩn, giao tiếp, IoT và giao diện, tổ chức lưu dữ liệu, tổng hợp chi phí. Sản phẩm phụ trách: hồ sơ đo lường, giao diện, nhật ký dữ liệu và bản báo cáo hợp nhất.
\end{itemize}

Ba thành viên cùng thống nhất thông số đầu vào, lắp ghép và đánh giá toàn hệ thống. Với phần cứng BMS, Lê Thanh Phú phối hợp Hoàng Ngô Thiên Phúc về đo lường; với đường công suất, Phan Trần Nguyên Phúc phối hợp Lê Thanh Phú về quyền sạc/xả. Mỗi phần trước khi tích hợp cần một thành viên khác kiểm tra chéo.

\subsection{Phối hợp và quản lý thay đổi}

Nhóm dự kiến rà tiến độ mỗi tuần, cập nhật công việc đã hoàn thành, vướng mắc và đầu ra tuần kế tiếp. Nhóm sẽ dùng chung bảng thông số gồm quy ước dòng, đơn vị, cấu hình pin, dải đo, ngưỡng và giao tiếp giữa các khối.

Khi thay loại pin, số cell, dải công suất hoặc cách chuyển nguồn, nhóm phải cập nhật mô hình, firmware, BOM và ca kiểm thử liên quan trước khi tiếp tục tích hợp. Mỗi mốc bàn giao sẽ lưu phiên bản tài liệu và mô hình tương ứng để truy lại dữ liệu.

\section{Các thông số cần xác định và rủi ro thực hiện}

\subsection{Các quyết định còn mở}

\begin{longtable}{@{}p{5.0cm} p{5.2cm} p{4.1cm}@{}}
\toprule
\textbf{Nội dung} & \textbf{Trạng thái hiện tại} & \textbf{Thời điểm cần xác định} \\
\midrule
\endfirsthead
\toprule
\textbf{Nội dung} & \textbf{Trạng thái hiện tại} & \textbf{Thời điểm cần xác định} \\
\midrule
\endhead
Công suất tải và thời gian hoạt động & Chưa xác định & M2, trước khi chốt quy mô nguồn và pin \\
Loại cell, số cell và dung lượng & Đề xuất khảo sát LiFePO4 4S; 15~Ah chỉ là giả định ban đầu & M2, trước thiết kế BMS chi tiết và mua pin \\
Ngưỡng điện áp/dòng/nhiệt, trễ và khôi phục & Chưa có bảng datasheet của pin được chọn & M2--M3, trước khi kết luận các ca bảo vệ \\
Dải đo và sai số cho phép & Chưa xác định từ phần cứng & M2; kiểm tra lại sau hiệu chuẩn ở M5 \\
Tham số PV thực tế & SM55 chỉ là bộ tham số tham khảo & M4, khi chọn tấm pin triển khai \\
MCU, cảm biến, thư viện và giấy phép Proteus & Chưa kiểm tra khả năng mô phỏng trên máy & Đầu M2 \\
Bộ sạc, inverter, nguồn phụ và ngân sách & Chưa chọn linh kiện hoặc mô-đun & M4 \\
Nền tảng IoT và giao thức điều khiển & ESP32/ThingSpeak là phương án khảo sát & M4--M5 \\
Thực hiện reconfiguration ngoài phần lý thuyết & Chưa đưa vào đầu ra bắt buộc & Xem xét sau khi hoàn thành các phần chính \\
\bottomrule
\caption{Các quyết định còn mở và thời điểm cần xác định.}
\label{tab:open-decisions}
\end{longtable}

\subsection{Rủi ro và cách xử lý dự kiến}

\begin{longtable}{@{}p{4.6cm} p{4.3cm} p{5.4cm}@{}}
\toprule
\textbf{Rủi ro} & \textbf{Ảnh hưởng} & \textbf{Cách xử lý} \\
\midrule
\endfirsthead
\toprule
\textbf{Rủi ro} & \textbf{Ảnh hưởng} & \textbf{Cách xử lý} \\
\midrule
\endhead
Thiếu mô hình MCU/cảm biến hoặc không nạp/lưu được dự án Proteus & Không chạy được một phần mô phỏng & Kiểm tra ngay đầu M2; dùng nguồn tín hiệu và khối cơ bản cho phần đo; giải quyết môi trường MCU trước kiểm thử SoC \\
Mô hình cell quá đơn giản & Dễ diễn giải sai thời gian sạc/xả hoặc cân bằng & Tách kiểm thử logic và kiểm thử động học; ghi rõ giới hạn của từng mô hình \\
Sai số đo lớn so với ngưỡng cân bằng & Bật sai nhánh hoặc rung đóng/cắt & Lập ngân sách sai số, hiệu chuẩn, điều chỉnh dải đo và ngưỡng \\
Mạch công suất không tương thích & Khó tích hợp hoặc vượt khả năng linh kiện & Rà điện áp, dòng, chiều cấp nguồn và điều kiện bảo vệ ở M4 \\
Linh kiện giao chậm hoặc phạm vi quá rộng & Thiếu thời gian thực nghiệm & Khảo sát cung ứng sớm; ưu tiên đầu ra cốt lõi, giữ tuần dự phòng và dùng mô-đun phù hợp khi cần \\
Mất mạng hoặc lệnh từ xa xung đột với bảo vệ & Hiển thị sai trạng thái hoặc điều khiển không phù hợp & Bảo vệ cục bộ, đánh dấu dữ liệu cũ, xác nhận lệnh và thử lỗi mạng \\
\bottomrule
\caption{Rủi ro thực hiện và cách xử lý dự kiến.}
\label{tab:risks}
\end{longtable}

\section{Sản phẩm dự kiến và phương pháp đánh giá}

\subsection{Nội dung đã trình bày trong báo cáo}

Báo cáo trình bày cơ sở lý thuyết, phương pháp thiết kế BMS 12~V, kế hoạch mô phỏng BMS/PV trên Proteus, phần lý thuyết PV array configuration, kiến trúc hệ thống và kế hoạch phân công. \Cref{tab:req-map} đối chiếu các yêu cầu của đề tài \cite{qa1} với mục tương ứng trong báo cáo.

\begin{longtable}{@{}p{6.5cm} p{8.0cm}@{}}
\toprule
\textbf{Yêu cầu} & \textbf{Mục trình bày trong báo cáo} \\
\midrule
\endfirsthead
\toprule
\textbf{Yêu cầu} & \textbf{Mục trình bày trong báo cáo} \\
\midrule
\endhead
BMS 12~V: SoC, nhiệt độ, cân bằng cell, bảo vệ & Mục 4.5--4.6 và Mục 5 \\
Mô phỏng chức năng trên Proteus trước phần cứng & Mục 6 và thứ tự M2--M5 trong Mục 9 \\
Mô phỏng tấm pin & Mục 4.2--4.3 và 6.2 \\
Tìm hiểu PV array configuration & Mục 7 \\
Phương pháp tổng quan, kế hoạch, phân nhiệm vụ & Mục 5--6 và 8--9 \\
Đọc tài liệu được cung cấp và tự tìm thêm & Mục 3 và 12 \\
\bottomrule
\caption{Đối chiếu yêu cầu của đề tài và nội dung trình bày trong báo cáo.}
\label{tab:req-map}
\end{longtable}

\subsection{Sản phẩm dự kiến của toàn project}

Các đầu ra dự kiến gồm mô hình Proteus PV và BMS; cấu hình và chương trình điều khiển; mô hình microgrid phần cứng với bộ sạc, pin, inverter, tải và chuyển nguồn; hệ thống IoT/giao diện; hồ sơ đo lường, thực nghiệm, BOM và báo cáo tổng kết.

Với các thành phần dùng mô-đun có sẵn, nhóm nêu mã sản phẩm, chức năng, cách tích hợp và phạm vi thử nghiệm. Phần tự phát triển có tài liệu đủ để tái lập mô hình và kiểm tra kết quả.

\subsection{Các chỉ tiêu đánh giá}

\begin{longtable}{@{}p{3.4cm} p{11.0cm}@{}}
\toprule
\textbf{Đối tượng} & \textbf{Chỉ tiêu và dữ liệu cần có} \\
\midrule
\endfirsthead
\toprule
\textbf{Đối tượng} & \textbf{Chỉ tiêu và dữ liệu cần có} \\
\midrule
\endhead
PV mô phỏng & Sai lệch các điểm đặc trưng so với nguồn tham chiếu; chiều biến thiên khi thay đổi \(G,T\) \\
Đo lường BMS & Sai số điện áp, dòng, nhiệt; đúng dấu dòng và đúng điện áp từng cell \\
SoC & Sai số so với điện tích tham chiếu trong ca thử có điều kiện ban đầu xác định; độ trôi khi có sai số dòng \\
Cân bằng & Chọn đúng nhánh, dòng cân bằng, chênh lệch cell trước/sau và thời gian theo mô hình được sử dụng \\
Bảo vệ & Điều kiện và thời gian tác động; dòng còn lại sau ngắt; điều kiện khôi phục; phản ứng với lỗi kết hợp \\
Nguồn và tải & Điện áp bus, khả năng cấp tải, quá trình chuyển nguồn và ảnh hưởng đến tải \\
IoT & Tính đúng của dữ liệu, độ trễ cập nhật/lệnh, phản hồi thực thi và hành vi khi mất mạng \\
Chi phí & Tổng BOM và chi phí thực hiện; so sánh các phương án đáp ứng cùng yêu cầu \\
\bottomrule
\caption{Các chỉ tiêu đánh giá cho từng đối tượng.}
\label{tab:metrics}
\end{longtable}

Mức chấp nhận cho từng chỉ tiêu được xác định trước khi chạy bộ thử chính thức, dựa trên yêu cầu tải, độ phân giải đo và linh kiện đã chọn. \pending{mức chấp nhận cụ thể cho từng chỉ tiêu}

Với giá trị tham chiếu khác 0, sai lệch tương đối được tính:
\[
\varepsilon_X=
\frac{|X_{\mathrm{thu}}-X_{\mathrm{ref}}|}{|X_{\mathrm{ref}}|}\times100\%.
\]

Với giá trị gần 0, nhóm báo cáo sai số tuyệt đối để tránh tỷ lệ không có ý nghĩa. Năng lượng ở một nhánh đo được tích phân từ công suất:
\[
E=\sum_k P_k\,\frac{\Delta t_k}{3600}\quad[\mathrm{Wh}],
\]
với công suất tính bằng W và thời gian bằng giây. Khi đánh giá hiệu suất, nhóm ghi rõ ranh giới khối và khoảng đo. Với hệ có pin, nhóm tính phần năng lượng lưu trữ thay đổi hoặc so sánh ở trạng thái pin đầu/cuối tương đương; không lấy điện năng tải chia cho điện năng PV trong một khoảng pin đang xả rồi kết luận đó là hiệu suất hệ thống.

Các thử nghiệm cuối gồm thay đổi mức tải, thay đổi mức nguồn PV hoặc điều kiện chiếu sáng ghi nhận được, chuyển nguồn phụ, mất kết nối và phục hồi sau lỗi. Nếu không đo được bức xạ, nhóm mô tả điều kiện thử thực tế thay vì gán giá trị \(\mathrm{W/m^2}\) tùy ý cho thời tiết.

% ============================================================
\section*{Tài liệu tham khảo}
\addcontentsline{toc}{section}{Tài liệu tham khảo}
% ============================================================

\begin{thebibliography}{9}

\bibitem{yaqoob}
S. J. Yaqoob, S. Motahhir và E. B. Agyekum,
``A new model for a photovoltaic panel using Proteus software tool under arbitrary environmental conditions,''
\textit{Journal of Cleaner Production}, tập 333, bài 130074, 2022.
DOI: \href{https://doi.org/10.1016/j.jclepro.2021.130074}{10.1016/j.jclepro.2021.130074}. Tài liệu do giảng viên cung cấp.

\bibitem{kulkarni}
P. Kulkarni, L. S. Paragond và S. Hiremath,
``Hybrid battery management system using the internet of things,''
\textit{Majlesi Journal of Electrical Engineering}, tập 19, số 2, bài 192531, tr. 1--8, 2025.
DOI: \href{https://doi.org/10.57647/j.mjee.2025.1902.31}{10.57647/j.mjee.2025.1902.31}. Tài liệu do giảng viên cung cấp.

\bibitem{thanh}
Ngô Ngọc Thành và Nguyễn Phùng Quang,
``Simulation of reconfiguration system using MATLAB--Simulink environment,''
\textit{Journal of Computer Science and Cybernetics}, tập 34, số 2, tr. 127--143, 2018.
DOI: \href{https://doi.org/10.15625/1813-9663/34/2/9194}{10.15625/1813-9663/34/2/9194}. Tài liệu do giảng viên cung cấp.

\bibitem{pvpmc}
M. Boyd và C. Hansen,
``Single Diode Equivalent Circuit Models,''
\textit{PV Performance Modeling Collaborative}, Sandia National Laboratories; nội dung đóng góp bởi NIST và Sandia.
\url{https://pvpmc.sandia.gov/modeling-guide/2-dc-module-iv/single-diode-equivalent-circuit-models/}. Truy cập ngày 10/09/2026.

\bibitem{tigauge}
N. Richards,
\textit{Battery Gauging Algorithm Comparison},
Texas Instruments, Application Note SLUAAR3, 12/2023.
\url{https://www.ti.com/lit/an/sluaar3/sluaar3.pdf}. Truy cập ngày 10/09/2026.

\bibitem{tibalance}
M. Sunna,
\textit{Cell Balancing With BQ769x2 Battery Monitors},
Texas Instruments, Application Note SLUAA81A, bản sửa đổi 02/2022.
\url{https://www.ti.com/lit/sluaa81}. Truy cập ngày 10/09/2026.

\bibitem{lm35}
Texas Instruments,
\textit{LM35 Precision Centigrade Temperature Sensors},
datasheet SNIS159H, bản sửa đổi 12/2017.
\url{https://www.ti.com/lit/ds/symlink/lm35.pdf}. Truy cập ngày 10/09/2026.

\bibitem{specs}
\textit{Chi tiết đề tài HK261-DAGD1-133: Thiết kế và phát triển hệ thống microgrid sử dụng năng lượng mặt trời với chi phí tối ưu kết hợp nền tảng IoT},
mô tả và yêu cầu học phần HK261 (bản lưu nội bộ: \texttt{docs/specifications/specs.md}).

\bibitem{qa1}
\textit{Q\&A với giảng viên, vòng 1 về proposal GĐ1},
trao đổi ngày 28/08/2026 (bản lưu nội bộ: \texttt{docs/info/qa1.md}).

\end{thebibliography}


\end{document}
